%% Abstract
\begin{abstract}

% %% Text of abstract
% As per \cite{alzint2025}, someone in the world develops dementia every 3 seconds. As of 2019, 55 million people worldwide suffered from dementia, and this number is expected to reach 139 million by 2050 \cite{who2023}. Alzheimer's disease is a progressive neurodegenerative disease that impairs cognitive function and it is the most common form of dementia and may contribute to 60–70\% of all cases \cite{who2025}. Upto 40\% of dementia cases could be prevented or delayed with early detection \cite{livingston2020}.

% \newline

% Longitudinal data enables the study of rates of change and temporal ordering of biomarker changes relative to the onset of dementia.
% Cross sectional data, on the other hand are influenced by intersubject variability. Longitudinal data attenuates this and therefore provide more accurate information regarding the onset and course of changes over time. \cite{mcdade2018longitudinal}

% \newline

% However, longitudinal datasets such as ADNI suffer from missing data problem due to patients not appearing for certain visits \cite{gabel2022retaining}. Generating missing longitudinal data in medical imaging datasets is crucial for enhancing temporal continuity, improving disease progression modelling, and supporting robust AI-driven diagnostics.  

% \newline

% We present SYNBADD, a novel method for missing data imputation using diffeomorphic deformation fields. Our model is conditioned on Age and cognitive disease labels, factors known to influence brain atrophy. By leveraging longitudinal data, our method models subject-specific aging and disease trajectories while improving time and memory efficiency by predicting deformations instead of entire
% brain volumes. Using deformation fields on real data for synthesis, also avoids problems related to distortion of a patient’s brain morphology when the scans are generated from scratch.

% Using the ADNI dataset, we demonstrate XX\%, XY\%, XZ\% improvements in SSIM, PSNR and MSE for missing scan synthesis over current SoTA methods. Additionally, we utilize an LDDMM loss based on voxelmorph to analyze the generated deformation XW\% improvement
% of our model over using only cross-sectional data.
\end{abstract}